\documentclass[a4paper,12pt,openany]{article}


\usepackage{parskip}
\usepackage{graphicx} % Required for inserting images



\usepackage[pdftex]{graphicx}

\usepackage[utf8]{inputenc}
\usepackage[T1]{fontenc}
\usepackage[british]{babel}

\usepackage{amsmath}
\usepackage{amsthm}

\usepackage{typearea}

\usepackage[toc,page]{appendix}
\usepackage{float}
% Acronyms
\usepackage{arydshln}





% Divide your \newtheorem commands into groups
% Comment out the necessary commands

\theoremstyle{plain}
\newtheorem{theorem}{Theorem}
%\newtheorem{lemma}{Lemma}
\newtheorem{corollary}{Corollary}
\newtheorem{proposition}{Proposition}

\theoremstyle{definition}
\newtheorem{definition}{Definition}
%\newtheorem{condition}{Condition}
%\newtheorem{problem}{Problem}
%\newtheorem{example}{Example}

\theoremstyle{remark}
%\newtheorem{remark}{Remark}
%\newtheorem{note}{Note}
%\newtheorem{claim}{Claim}

\usepackage{newtxtext}
\usepackage{newtxmath}


\usepackage{amsmath}

\begin{document}


\begin{center}
\huge Algebraic Geometry\\
\vspace{0.2cm}
\Large Solutions to selected problems for assignment 2 \\  \vspace{0.5cm}
Mara Kali\v{c}anin Dimitrov
\vspace{0.5cm}
\end{center}

\subsection*{Problem 7.2}
Let $k$ be an algebraically closed field of characteristic $0$,
and let $C\subseteq \mathbb{A}^2$ be the affine plane curve defined by
\[
f(x,y)=y^{2}-x^{3}-x^{2}\in k[x,y].
\]
We need to:
\begin{itemize}
\item[i)] show that $P=(0,0)$ is the \emph{only} multiple (singular) point of $C$,
\item[ii)] and find the multiplicity of $C$ at $P$.
\end{itemize}

\medskip

Let's recall the relevant definitions from Olof's book:
\begin{itemize}
\item  \textbf{Definition 1.} If $C=V(f)\subseteq \mathbb{A}^2$ and $Q\in C$, then $Q$ is \emph{nonsingular (simple)} if
\[
\frac{\partial f}{\partial x}(Q)\neq 0 \quad \text{or}\quad \frac{\partial f}{\partial y}(Q)\neq 0,
\]
and $Q$ is \emph{singular (multiple)} if both partial derivatives vanish at $Q$.

\item \textbf{Definition 7.11.} For any $f\in k[x,y]$ can be written uniquely as
\[
f=f_m+f_{m+1}+\cdots+f_d
\]
where each $f_i$ is homogeneous of degree $i$ and $f_m\neq 0$.

\item \textbf{Definition 7.12.} For $C=V(f)$, the multiplicity $m_{(0,0)}(C)$ equals the least $m$
such that the homogeneous part $f_m$ is nonzero.
\end{itemize}

\medskip

By Definition 1, a point $Q=(a,b)\in C$ is singular if and only if it satisfies the \emph{system}
\[
f(a,b)=0,\qquad \frac{\partial f}{\partial x}(a,b)=0,\qquad \frac{\partial f}{\partial y}(a,b)=0.
\]

The partial derivatives are 
\[
\frac{\partial f}{\partial x}(x,y)=\frac{\partial}{\partial x}(y^2-x^3-x^2)=-x(3x+2),
\qquad
\frac{\partial f}{\partial y}(x,y)=2y.
\]

Let $Q=(a,b)\in C$ be singular. Then $\frac{\partial f}{\partial y}(a,b)=0$ gives
\[
2b=0.
\]
Since $\operatorname{char}(k)=0$, we have $2\neq 0$ hence $b=0$.

Now impose $Q\in C$, i.e.
\[
0=f(a,0)=-a^{2}(a+1),
\]
so
\[
a=0 \quad\text{or}\quad a=-1.
\]

Finally, impose $\frac{\partial f}{\partial x}(a,0)=0$:
\[
0=\frac{\partial f}{\partial x}(a,0)=-a(3a+2),
\]
so
\[
a=0 \quad\text{or}\quad a=-\frac{2}{3}.
\]

Thus $a$ must be $0$.
Therefore, we can conclude that $a=0$ and $b=0$, and the only singular point is
\[
P=(0,0).
\]
So we can conclude that $P=(0,0)$ is the unique multiple point of $C$.

\medskip


Now we compute $m_{P}(C)=m_{(0,0)}(C)$ using Definition 7.12.

Write $f$ as a sum of homogeneous parts as in Definition 7.11
\[
f(x,y)=\underbrace{(y^{2}-x^{2})}_{=:f_{2}(x,y)} \;+\; \underbrace{(-x^{3})}_{=:f_{3}(x,y)}.
\]
Here $f_2$ is homogeneous of degree $2$, $f_3$ is homogeneous of degree $3$, and $f_2\neq 0$.
Hence, the lowest nonzero homogeneous part occurs in degree $2$, so by Definition 7.12 we can conclude that
\[
m_{(0,0)}(C)=2.
\]



\subsection*{Problem 7.3}

Let $k$ be the algebraically closed field of characteristic $0$, and let
\[
f(x,y)=y^{2}-x^{3}\in k[x,y],
\qquad
C:= V(f)\subseteq \mathbb{A}^{2}.
\]
We need to:
\begin{itemize}
\item[i)]  show $P=(0,0)$ is the \emph{only} multiple (singular) point of $C$;
\item[ii)] find the multiplicity of $C$ at $P$.
\end{itemize}

\medskip

By Definition 1, a point $Q=(a,b)\in C$ is singular if and only if \emph{both} partial derivatives vanish at $Q$.
Equivalently, to find singular points we must solve the system
\[
f(a,b)=0,\qquad \frac{\partial f}{\partial x}(a,b)=0,\qquad \frac{\partial f}{\partial y}(a,b)=0.
\]
The partial derivatives are

\[
\frac{\partial f}{\partial x}(x,y)=-3x^{2},
\qquad
\frac{\partial f}{\partial y}(x,y)=2y.
\]
Let $Q=(a,b)\in C$ be singular. Then
\[
-3a^{2}=0
\quad\text{and}\quad
2b=0.
\]
Since $\operatorname{char}(k)=0$, we have $2,3\neq 0$, so these equations are equivalent to
\[
a^{2}=0
\quad\text{and}\quad
b=0.
\]
Because $k$ is a field, $a^{2}=0$ implies $a=0$. Hence, the only possible singular point is $Q=(0,0)$.

Finally, we check that $(0,0)$ lies on the curve
\[
f(0,0)=0^{2}-0^{3}=0,
\]
so $(0,0)\in V(f)=C$.

Therefore $P=(0,0)$ is the \emph{only} multiple (singular) point of $C$.

\medskip

By Definition 7.12, $m_{(0,0)}(C)$ is the degree of the \emph{lowest nonzero homogeneous part} of $f$.
So we write $f$ as a sum of homogeneous terms and read off the smallest degree that occurs.

The polynomial $y^{2}$ is homogeneous of degree $2$ and $-x^{3}$ is homogeneous of degree $3$, hence
\[
f(x,y)=y^{2}-x^{3}= f_{2}+f_{3},
\qquad
f_{2}=y^{2}\neq 0,\quad f_{3}=-x^{3}.
\]
Thus, the lowest degree homogeneous part is $f_{2}$, which has degree $2$. Therefore by Definition 7.12, we can conclude that
\[
m_{(0,0)}(C)=2.
\]
Hence $m_{P}(C)=2$.


\subsection*{Problem 7.4}
Let $k$ be an algebraically closed field of characteristic $0$, and let
$C\subseteq \mathbb{A}^2$ be the affine plane curve defined by
\[
f(x,y)=y^{3}-y^{2}+x^{3}-x^{2}+3x^{2}y+3xy^{2}+2xy
\in k[x,y].
\]
We need to do the following
\begin{itemize}
\item[a)] determine all \emph{singular (multiple)} points of $C$;
\item[b)] for each singular point $P$, determine the \emph{tangent line(s)} of $C$ at $P$.
\end{itemize}

\medskip

Let's recall the relevant definition from Olof's book:
\begin{itemize}
    \item \textbf{Definition 7.15}.  if $f_m$ is the lowest-degree homogeneous part and
    \[
    f_m=\prod_{i=1}^r g_i^{\,r_i}
    \]
    is the factorization of $f_m$ into \emph{distinct} linear factors $g_i$, then the lines $L_i=V(g_i)$ are the tangent lines
    of $C$ at $(0,0)$, with multiplicities $r_i$.
\end{itemize}

\medskip
\textbf{a)} By Definition 1, a point $P=(a,b)\in C$ is singular if and only if both partial derivatives vanish at $P$.
Equivalently, to \emph{find} the singular points of $C$ we solve the system
\[
f(a,b)=0,\qquad \frac{\partial f}{\partial x}(a,b)=0,\qquad \frac{\partial f}{\partial y}(a,b)=0.
\]
The partial derivatives are given by
\[
\frac{\partial f}{\partial x}(x,y)=3x^{2}-2x+6xy+3y^{2}+2y,
\qquad
\frac{\partial f}{\partial y}(x,y)=3y^{2}-2y+3x^{2}+6xy+2x.
\]
Let $(a,b)\in \mathbb{A}^2$ satisfy $\frac{\partial f}{\partial x}(a,b)=\frac{\partial f}{\partial y}(a,b)=0$.
Subtracting the two equations gives
\begin{align*}
0
&=\frac{\partial f}{\partial x}(a,b)-\frac{\partial f}{\partial y}(a,b)\\
&=\bigl(3a^{2}-2a+6ab+3b^{2}+2b\bigr)-\bigl(3b^{2}-2b+3a^{2}+6ab+2a\bigr)\\
&=4(b-a).
\end{align*}
Since $4\neq 0$, we conclude $b=a$.

Substitute $b=a$ into $\frac{\partial f}{\partial y}(a,b)=0$ we get the following
\[
0=\frac{\partial f}{\partial y}(a,a)=3a^{2}-2a+3a^{2}+6a^{2}+2a=12a^{2}.
\]
Since $12\neq 0$, it follows that $a^{2}=0$, hence $a=0$ (because $k$ is a field), and therefore $b=a=0$.

Thus the only \emph{candidate} point where both partial derivatives vanish is $(0,0)$.
We now check that it lies on the curve:
\[
f(0,0)=0^{3}-0^{2}+0^{3}-0^{2}+0=0,
\]
so $(0,0)\in C$.

Therefore, $P=(0,0)$ is a singular point of $C$, and it is the \emph{only} singular point.

\medskip

\textbf{b)} Since part a) shows that the only singular point is $P=(0,0)$, we may apply Definition 7.15 directly at the origin.

The degree $2$ part is
\[
f_{2}(x,y)= -y^{2}-x^{2}+2xy = -(x^{2}-2xy+y^{2})=-(x-y)^{2},
\]
and the remaining terms have degree $3$:
\[
f_{3}(x,y)=y^{3}+x^{3}+3x^{2}y+3xy^{2}.
\]
Thus $f=f_{2}+f_{3}$ and the lowest-degree homogeneous part is $f_{2}$.

Since $-1\neq 0$, the factorization of $f_{2}$ into distinct linear factors is
\[
f_{2}(x,y)=-(x-y)^{2}=(-1)\cdot (x-y)^{2}.
\]
Therefore, by Definition 7.15, $C$ has exactly one tangent line at $(0,0)$, namely
\[
L:\ x-y=0 \quad (\text{i.e.\ } y=x),
\]
and this tangent line occurs with multiplicity $2$.

\medskip





\subsection*{Problem 7.5}

Let $k$ be an algebraically closed field of characteristic $0$, and let
$C\subseteq \mathbb{A}^2$ be the affine plane curve defined by
\[
f(x,y)=x^4+y^4-x^2y^2 \in k[x,y].
\]

We need to do the following
\begin{itemize}
\item[a)] determine all \emph{singular (multiple)} points of $C$;
\item[b)] for each singular point $P$, determine the \emph{tangent line(s)} of $C$ at $P$.
\end{itemize}

\medskip

\textbf{a)}
By Definition 1, a point $P=(a,b)\in C$ is singular if and only if it satisfies the system
\[
f(a,b)=0,\qquad \frac{\partial f}{\partial x}(a,b)=0,\qquad \frac{\partial f}{\partial y}(a,b)=0.
\]

First, compute the partial derivatives
\[
\frac{\partial f}{\partial x}(x,y)=2x(2x^{2}-y^{2}),
\qquad
\frac{\partial f}{\partial y}(x,y)=2y(2y^{2}-x^{2}).
\]
Let $P=(a,b)\in C$ be singular. Since $2\neq 0$ then we have that
\[
a(2a^{2}-b^{2})=0,
\qquad
b(2b^{2}-a^{2})=0.
\]
Thus
\[
a=0 \ \text{or}\ b^{2}=2a^{2},
\qquad
b=0 \ \text{or}\ a^{2}=2b^{2}.
\]

Therefore, we have the following cases

\emph{Case 1: $a=0$.} Then $f(0,b)=b^{4}=0$, hence $b=0$. Thus $P=(0,0)$.

\emph{Case 2: $b=0$.} Then $f(a,0)=a^{4}=0$, hence $a=0$. Thus $P=(0,0)$.

\emph{Case 3: $a\neq 0$ and $b\neq 0$.} Then we must have
\[
b^{2}=2a^{2}\quad\text{and}\quad a^{2}=2b^{2}.
\]
Substituting the first into the second gives
\[
a^{2}=2(2a^{2})=4a^{2}\quad\Longrightarrow\quad 3a^{2}=0.
\]
Since $\operatorname{char}(k)=0$, we have $3\neq 0$, hence $a^{2}=0$ and therefore $a=0$,
contradicting $a\neq 0$. Hence this case is impossible.

Therefore, the only singular point of $C$ is
\[
P=(0,0).
\]

\medskip

\textbf{b)} By Definition 7.15, the tangent lines at $(0,0)$ are obtained from the factorization of the \emph{lowest-degree homogeneous part} $f_m$.
Here $f$ is already homogeneous of degree $4$, so by Definition 7.11 we have
\[
f=f_4,
\]
and thus $f_m=f_4=f$. 

By Definition 7.12  we can conclude that $m_{(0,0)}(C)=4$.

So it remains to factor the homogeneous form $f$ into linear factors.
Since $k$ is algebraically closed, choose $\sqrt{3}\in k$ with $(\sqrt{3})^{2}=3$. Then
\[
(x^{2}+y^{2}+\sqrt{3}xy)(x^{2}+y^{2}-\sqrt{3}xy)
=(x^{2}+y^{2})^{2}-(\sqrt{3}xy)^{2}
=x^{4}+y^{4}-x^{2}y^{2},
\]
so
\[
f=(x^{2}+y^{2}+\sqrt{3}xy)(x^{2}+y^{2}-\sqrt{3}xy).
\]

Next, again since $k$ is algebraically closed, choose $i\in k$ with $i^{2}=-1$.
Consider the first quadratic as a polynomial in $y$:
\[
x^{2}+y^{2}+\sqrt{3}xy = y^{2}+(\sqrt{3}x)y+x^{2}.
\]
Hence the roots are
\[
y=\frac{-\sqrt{3}x\pm ix}{2}=\left(\frac{-\sqrt{3}\pm i}{2}\right)x,
\]
and therefore
\[
x^{2}+y^{2}+\sqrt{3}xy
=
\left(y-\frac{-\sqrt{3}+i}{2}x\right)\left(y-\frac{-\sqrt{3}-i}{2}x\right).
\]
Similarly,
\[
x^{2}+y^{2}-\sqrt{3}xy
=
\left(y-\frac{\sqrt{3}+i}{2}x\right)\left(y-\frac{\sqrt{3}-i}{2}x\right).
\]

Putting these together, we obtain a factorization of $f$ into \emph{distinct} linear factors:
\[
f
=
\left(y-\frac{-\sqrt{3}+i}{2}x\right)
\left(y-\frac{-\sqrt{3}-i}{2}x\right)
\left(y-\frac{\sqrt{3}+i}{2}x\right)
\left(y-\frac{\sqrt{3}-i}{2}x\right).
\]
Therefore, by Definition 7.15, the tangent lines of $C$ at $P=(0,0)$ are exactly the four lines
\[
L_1:\ y-\frac{-\sqrt{3}+i}{2}x=0,\quad
L_2:\ y-\frac{-\sqrt{3}-i}{2}x=0,\quad
L_3:\ y-\frac{\sqrt{3}+i}{2}x=0,\quad
L_4:\ y-\frac{\sqrt{3}-i}{2}x=0.
\]
Each linear factor appears to the first power, so each tangent line has multiplicity $1$.

\medskip

\subsection*{Problem 7.6}

Let $C\subseteq \mathbb{A}^{2}$ be the curve defined by
\[
f(x,y)=(x^{2}+y^{2}-3x)^{2}-4x^{2}(2-x)\in k[x,y],
\]
where $k$ is algebraically closed of characteristic $0$.

\begin{itemize}
\item[(a)] Compute the multiplicity $m_{P}(C)$ at $P=(0,0)$.
\item[(b)] Find the other singular point $Q$ of $C$ and compute its multiplicity $m_{Q}(C)$.
\end{itemize}

\medskip

Let's recall the relevant definition and result from Olof's book:
\begin{itemize}
\item \textbf{Definition 7.13.} Let $C$ be an affine plane curve defined by the polynomial $f$ and let $P$ be any point. Let $\phi$ be change of coordinates of $\mathbb{A}^2$ taking $P$ to $Q=(0,0)$ and let $D$ be the curve defined by $g=\phi^* f$. The multiplicity $m_P(C)$ of $C$ at $P$ is then defined as $m_Q(D)$.

\item \textbf{Corollary 6.41}. Two varieties $X$ and $Y$ are isomorphic if and only if their coordinate rings $\Gamma(X)$ and $\Gamma(Y)$ are isomorphic.

An isomorphism $\phi: \mathbb{A}^n \rightarrow \mathbb{A}^n$ such that each $f_i$ has degree 1 is also called a change of coordinates. If we view $\mathbb{A}^n$ as a vector space, we can factor $\phi$ as $\phi=\tau \circ \lambda$ where $\tau$ is a translation and $\lambda$ is an invertible linear map. Conversely, any composition of translations and invertible linear maps give a change of coordinates.
\end{itemize}

\medskip


\textbf{a)} By Definition 7.12, $m_{(0,0)}(C)$ is the least total degree of a nonzero homogeneous part of $f$. First note that $P\in C$, since
\[
f(0,0)=(0+0-0)^{2}-0=0.
\]

Now expand $f$ as
\[
(x^{2}+y^{2}-3x)^{2}=(x^{2}+y^{2})^{2}-6x(x^{2}+y^{2})+9x^{2},
\qquad
4x^{2}(2-x)=8x^{2}-4x^{3}.
\]
Hence
\begin{align*}
f(x,y)
&=(x^{2}+y^{2})^{2}-6x(x^{2}+y^{2})+9x^{2}-(8x^{2}-4x^{3})\\
&=(x^{2}+y^{2})^{2}-6x(x^{2}+y^{2})+x^{2}+4x^{3}\\
&=x^{4}+2x^{2}y^{2}+y^{4}-2x^{3}-6xy^{2}+x^{2}.
\end{align*}
Grouping by total degree by Definition 7.11, we get
\[
f(x,y)=\underbrace{x^{2}}_{=:f_{2}}
\;+\;
\underbrace{\bigl(-2x^{3}-6xy^{2}\bigr)}_{=:f_{3}}
\;+\;
\underbrace{\bigl(x^{4}+2x^{2}y^{2}+y^{4}\bigr)}_{=:f_{4}}.
\]
The lowest-degree nonzero part is $f_{2}=x^{2}\neq 0$, and it has degree $2$. Therefore, by Definition 7.12,
\[
m_{P}(C)=m_{(0,0)}(C)=2.
\]


\medskip

\textbf{b)} 
By Definition 1, $R=(a,b)\in C$ is singular if and only if
\[
f(a,b)=0,\qquad \frac{\partial f}{\partial x}(a,b)=0,\qquad \frac{\partial f}{\partial y}(a,b)=0.
\]
The partial derivatives are
\[
\frac{\partial f}{\partial x}
=4x^{3}-6x^{2}+4xy^{2}+2x-6y^{2},
\qquad
\frac{\partial f}{\partial y}
=4y(x^{2}-3x+y^{2}).
\]
Since $4\neq 0$ we get
\[
y(x^{2}-3x+y^{2})=0.
\]
We consider the two cases coming from this product.

\emph{Case 1: $y=0$.}
Then
\[
f(x,0)=x^{4}-2x^{3}+x^{2}=x^{2}(x-1)^{2},
\]
so the points of $C$ with $y=0$ are $(0,0)$ and $(1,0)$.
For these points to be singular we also need $\frac{\partial f}{\partial x}(x,0)=0$. Compute
\[
\frac{\partial f}{\partial x}(x,0)
=2x(2x^{2}-3x+1)=2x(2x-1)(x-1).
\]
Thus $\frac{\partial f}{\partial x}(0,0)=0$ and $\frac{\partial f}{\partial x}(1,0)=0$.
Since also $\frac{\partial f}{\partial y}(x,0)=0$ for all $x$, both $(0,0)$ and $(1,0)$ are singular points of $C$.

\emph{Case 2: $x^{2}-3x+y^{2}=0$.}
Then $x^{2}+y^{2}-3x=0$. Using the original defining equation
\[
f(x,y)=(x^{2}+y^{2}-3x)^{2}-4x^{2}(2-x),
\]
this implies
\[
f(x,y)=0 \quad\Longleftrightarrow\quad -4x^{2}(2-x)=0
\quad\Longleftrightarrow\quad x^{2}(2-x)=0.
\]
Hence $x=0$ or $x=2$. If $x=0$, then $x^{2}-3x+y^{2}=0$ gives $y^{2}=0$, hence $y=0$, which is already covered in Case~1.
If $x=2$, then $x^{2}-3x+y^{2}=0$ gives $y^{2}=2$. For such points,
\[
\frac{\partial f}{\partial x}(2,y)
=4\cdot 2^{3}-6\cdot 2^{2}+4\cdot 2\cdot y^{2}+2\cdot 2-6y^{2}
=12+2y^{2}.
\]
If $y^{2}=2$, then $\frac{\partial f}{\partial x}(2,y)=12+4=16\neq 0$, so these points are \emph{not} singular.

Combining the two cases, we conclude that the singular points of $C$ are exactly
\[
P=(0,0)\quad\text{and}\quad Q=(1,0).
\]

\medskip

Next we want to compute $m_{Q}(C)$.
Multiplicity at a general point is defined by translating that point to the origin and applying Definition 7.12 which is our Definition 7.13. By Corollary 6.41, translations are changes of coordinates, so we may set
\[
u:=x-1,\qquad v:=y,
\]
so that $Q=(1,0)$ becomes $(u,v)=(0,0)$. In these coordinates the defining polynomial becomes
\[
g(u,v):=f(u+1,v).
\]
A direct expansion gives
\[
g(u,v)=u^{4}+2u^{3}+2u^{2}v^{2}+u^{2}-2uv^{2}+v^{4}-4v^{2}.
\]
Now write $g$ as a sum of homogeneous parts (Definition 7.11):
\[
g(u,v)=\underbrace{(u^{2}-4v^{2})}_{=:g_{2}}
+\underbrace{(2u^{3}-2uv^{2})}_{=:g_{3}}
+\underbrace{(u^{4}+2u^{2}v^{2}+v^{4})}_{=:g_{4}}.
\]
The lowest-degree nonzero homogeneous part is $g_{2}=u^{2}-4v^{2}\neq 0$, which has degree $2$.
Therefore, by Definition 7.12,
\[
m_{(0,0)}(V(g))=2,
\]
and hence by Definition 7.13,
\[
m_{Q}(C)=2.
\]

\medskip


\subsection*{Problem 7.10}

Let
\[
f(x,y)=y^{2}-x^{3}-x^{2},\qquad g(x,y)=y^{2}-x^{3}\in k[x,y],
\]
let $C=V(f)$ and $D=V(g)$ be the corresponding plane curves in $\mathbb{A}^{2}$, and let $P=(0,0)$.
We need to compute the local intersection number $i_{P}(C,D)$.

\medskip

Let's recall the relevant definition and results from Olof's book:
\begin{itemize}
    \item \textbf{Definition 6.3}. For $f,g\in k[x,y]$,
    \[
    V(f,g)=\{Q\in \mathbb{A}^2 \mid f(Q)=0 \text{ and } g(Q)=0\}.
    \]
    \item \textbf{Proposition 6.48}. Let $I$ be an ideal of $k\left[x_1, \ldots, x_n\right]$ such that $V(I)= \left\{P_1, \ldots, P_m\right\}$ is a finite set of points. Then

$$
k\left[x_1, \ldots, x_n\right] / I \cong \prod_{i=1}^m \mathcal{O}_{P_i}\left(\mathbb{A}^n\right) / I \mathcal{O}_{P_i}\left(\mathbb{A}^n\right)
$$
    \item \textbf{Theorem 7.29}. Let $C$ and $D$ be plane curves, defined by polynomials $f$ and $g$ respectively, and let $P \in \mathbb{A}^2$. The intersection number of $C$ and $D$ at $P$ is uniquely defined and given by

$$
i_P(C, D)=\operatorname{dim}_k\left(\mathcal{O}_P\left(\mathbb{A}^2\right) /(f, g)\right) .
$$
\end{itemize}

\medskip


By Theorem 7.29, computing $i_{P}(C,D)$ reduces to computing the $k$-dimension of the local ring quotient
$\mathcal{O}_{P}(\mathbb{A}^{2})/(f,g)$.  
To compute this dimension concretely, first we need to show that $V(f,g)=\{P\}$.
Let $(a,b)\in \mathbb{A}^{2}$ satisfy $f(a,b)=0$ and $g(a,b)=0$. Subtracting the two equations gives
\[
0=f(a,b)-g(a,b)=\bigl(b^{2}-a^{3}-a^{2}\bigr)-\bigl(b^{2}-a^{3}\bigr)=-a^{2}.
\]
Since $k$ is a field, $a^{2}=0$ implies $a=0$. Substituting $a=0$ into $g(a,b)=0$ yields
\[
0=g(0,b)=b^{2},
\]
so $b=0$. Hence $(a,b)=(0,0)=P$. Therefore
\[
V(f,g)=\{P\}.
\]

\medskip

Next set $I:=(f,g)\subseteq k[x,y]$. Since $V(I)=V(f,g)=\{P\}$, Proposition 6.48 (p.~88) gives an isomorphism
\[
k[x,y]/(f,g)\ \cong\ \mathcal{O}_{P}(\mathbb{A}^{2})/(f,g)\mathcal{O}_{P}(\mathbb{A}^{2}).
\]
Taking $k$-dimensions and applying Theorem 7.29, we obtain
\[
i_{P}(C,D)=\dim_{k}\!\bigl(\mathcal{O}_{P}(\mathbb{A}^{2})/(f,g)\bigr)
=\dim_{k}\!\bigl(k[x,y]/(f,g)\bigr).
\]
So it remains to compute the dimension of $k[x,y]/(f,g)$.

\medskip

Let's simplify $(f,g)$.The following is true
\[
g-f = (y^{2}-x^{3})-(y^{2}-x^{3}-x^{2})=x^{2}.
\]
Hence
\[
(f,g)=(f,\,g-f)=(f,\,x^{2}).
\]
We claim that
\[
(f,g)=(x^{2},y^{2}).
\]

First, $(x^{2},y^{2})\subseteq (f,g)$.
We already have $x^{2}\in(f,g)$, hence $x^{3}=x\cdot x^{2}\in(f,g)$.
Since $g=y^{2}-x^{3}$, we can rewrite
\[
y^{2}=g+x^{3}\in(f,g).
\]
Thus $x^{2},y^{2}\in(f,g)$, so $(x^{2},y^{2})\subseteq (f,g)$.

\emph{Conversely, $(f,g)\subseteq (x^{2},y^{2})$.}
Because $x^{3}=x\cdot x^{2}\in (x^{2})\subseteq (x^{2},y^{2})$, we have
\[
g=y^{2}-x^{3}\in (x^{2},y^{2})
\quad\text{and}\quad
f=y^{2}-x^{3}-x^{2}\in (x^{2},y^{2}).
\]
Therefore $(f,g)\subseteq (x^{2},y^{2})$.

Combining both inclusions gives
\[
(f,g)=(x^{2},y^{2}),
\]
hence
\[
k[x,y]/(f,g)\ \cong\ k[x,y]/(x^{2},y^{2}).
\]

\medskip

And finally, compute $\dim_{k}\bigl(k[x,y]/(x^{2},y^{2})\bigr)$.
In the quotient by $(x^{2},y^{2})$, any monomial $x^{a}y^{b}$ with $a\ge 2$ or $b\ge 2$ becomes $0$,
since it is divisible by $x^{2}$ or by $y^{2}$.
Therefore, every class has a representative of the form
\[
a+bx+cy+dxy \qquad (a,b,c,d\in k),
\]
so the images of the monomials $1,x,y,xy$ span the quotient as a $k$-vector space.

To see that these four classes are linearly independent, suppose that
\[
a+bx+cy+dxy \in (x^{2},y^{2}).
\]
Then there exist polynomials $P,Q\in k[x,y]$ such that
\[
a+bx+cy+dxy = x^{2}P + y^{2}Q.
\]
Every monomial on the right-hand side has $x$-exponent $\ge 2$ or $y$-exponent $\ge 2$,
while the left-hand side is a $k$-linear combination of $1,x,y,xy$, all of which have $x$- and $y$-exponents $\le 1$.
By the uniqueness of polynomial expressions in monomials, it follows that $a=b=c=d=0$.
Hence $\{1,\overline{x},\overline{y},\overline{x}\,\overline{y}\}$ is a $k$-basis and
\[
\dim_{k}\bigl(k[x,y]/(x^{2},y^{2})\bigr)=4.
\]

\medskip


Therefore
\[
i_{P}(C,D)=4.
\]


\subsection*{Problem 7.11}

Let
\[
f(x,y)=(y-x^{2})(y+x^{2}), \qquad g(x,y)=y^{2}-x^{3}-x^{2}\in k[x,y],
\]
let $C:=V(f)$ and $D:=V(g)$ be plane curves in $\mathbb{A}^{2}$, and let $P=(0,0)$.
We need to compute $i_{P}(C,D)$.

\medskip

Let's recall the relevant definitions and results from Olof's book:
\begin{itemize}
\item \textbf{Proposition 7.32}.
Suppose that $f=\prod_{i=1}^n f_i^{r_i}$ and that $g=\prod_{i=1}^m g_i^{s_i}$. Then $$i_P(f, g)=\sum_{i=1}^n \sum_{j=1}^m r_i s_j i_P\left(f_i, g_j\right).
$$
We use this because $f$ factors as $(y-x^2)(y+x^2)$, so $C$ has two components and the intersection number splits into a sum of two contributions.

\item \textbf{Theorem 7.23}. An irreducible affine plane curve $C$ is nonsigular at $P \in C$ if and only if $\mathcal{O}_{C, P}$ is a discrete valuation ring (DVR).

If $C$ is smooth at $P$, then a uniformizing parameter for $\mathcal{O}_{C, P}$ is given by $\bar{l}$, where $l \in k\left[x_1, x_2\right]$ is any polynomial defining a line $L$ through $P$ which is not tangent to $C$ at $P$.


\item \textbf{Proposition 7.34}. Let $C$ be a curve defined by the polynomial $f$, let $D$ be a curve defined by a polynomial $g$ and let $P$ be a nonsingular point of $C$. Then

$$
i_P(C, D)=\operatorname{ord}_P^C(g) .
$$

\item \textbf{Proposition 6.45}. The ring $\mathcal{O}_{X, P}$ is a local ring with maximal ideal
$$
\mathfrak{m}_P(X)=\left\{f \in \mathcal{O}_{X, P} \mid f(P)=0\right\} .
$$
\end{itemize}

\medskip

First, we will reduce to the two components of $C$.
Since $f$ factors as a product of two polynomials, the curve $C=V(f)$ is the union of two curves
\[
f=(y-x^{2})(y+x^{2})=:f_{+}f_{-},
\qquad
C_{+}:=V(f_{+}),\ \ C_{-}:=V(f_{-}),
\qquad
C=C_{+}\cup C_{-}.
\]
Instead of intersecting $D$ with the union $C$, we will intersect $D$ with each component and add the results.
By Proposition 7.32 we have that
\[
i_{P}(C,D)=i_{P}(f,g)=i_{P}(f_{+},g)+i_{P}(f_{-},g)
= i_{P}(C_{+},D)+i_{P}(C_{-},D).
\]
Thus it suffices to compute $i_{P}(C_{+},D)$ and $i_{P}(C_{-},D)$.

\medskip

Next, we want to check that $P$ is nonsingular on each component (so that Proposition 7.34 applies).
For $C_{+}$, the defining polynomial is $f_{+}(x,y)=y-x^{2}$. Then
\[
\frac{\partial f_{+}}{\partial y}=1,
\]
so $\frac{\partial f_{+}}{\partial y}(P)=1\neq 0$. Hence $P$ is a simple point of $C_{+}$ by Definition 1.

Similarly, for $C_{-}$ we have $f_{-}(x,y)=y+x^{2}$ and again
\[
\frac{\partial f_{-}}{\partial y}=1,
\]
so $\frac{\partial f_{-}}{\partial y}(P)=1\neq 0$. Hence $P$ is a simple point of $C_{-}$.

Therefore by Theorem 7.23 applies to both $C_{+}$ and $C_{-}$, and Proposition 7.34  will allow us to compute
intersection numbers by taking orders.

\medskip


Since $P$ is nonsingular on $C_{+}$, Proposition 7.34 gives
\[
i_{P}(C_{+},D)=\operatorname{ord}^{C_{+}}_{P}(g).
\]

To compute this order, we first restrict $g$ to the curve $C_{+}$.
In the coordinate ring $\Gamma(C_{+})=k[x,y]/(y-x^{2})$, we have $\bar{y}=\bar{x}^{2}$, so
\[
g|_{C_{+}} = g(x,x^{2}) = (x^{2})^{2}-x^{3}-x^{2}
= x^{4}-x^{3}-x^{2}
= x^{2}(x^{2}-x-1).
\]

Next, we choose a uniformizing parameter.
The tangent line to $C_{+}$ at $P$ is $y=0$ (Proved in Example 7.6. from Olof's book). In particular, the line $x=0$ is \emph{not} tangent to $C_{+}$ at $P$.
Thus, by Theorem 7.23, $x$ is a uniformizing parameter for the DVR $\mathcal{O}_{C_{+},P}$, so
\[
\operatorname{ord}^{C_{+}}_{P}(x)=1.
\]

Finally, note that
\[
(x^{2}-x-1)\big|_{x=0}=-1\neq 0,
\]
so $x^{2}-x-1$ does not vanish at $P$ on $C_{+}$ and is therefore a unit in $\mathcal{O}_{C_{+},P}$ by Proposition 6.45.
Hence it contributes order $0$, and we obtain
\[
\operatorname{ord}^{C_{+}}_{P}\bigl(x^{2}(x^{2}-x-1)\bigr)
=\operatorname{ord}^{C_{+}}_{P}(x^{2})
=2\operatorname{ord}^{C_{+}}_{P}(x)=2.
\]
Therefore
\[
i_{P}(C_{+},D)=2.
\]

\medskip

Next, compute $i_{P}(C_{-},D)$ in the same way.
Again Proposition 7.34 gives
\[
i_{P}(C_{-},D)=\operatorname{ord}^{C_{-}}_{P}(g).
\]
On $C_{-}$ we have $y=-x^{2}$, so
\[
g|_{C_{-}} = g(x,-x^{2}) = (-x^{2})^{2}-x^{3}-x^{2}
= x^{4}-x^{3}-x^{2}
= x^{2}(x^{2}-x-1).
\]
As before, $P$ is a simple point of $C_{-}$ and its tangent line at $P$ is again $y=0$, so $x=0$ is not tangent.
Thus $x$ is a uniformizing parameter in $\mathcal{O}_{C_{-},P}$ by Theorem 7.23, and the factor
$x^{2}-x-1$ is a unit because it evaluates to $-1$ at $x=0$ (Proposition 6.45). Hence
\[
\operatorname{ord}^{C_{-}}_{P}(g)=2,
\qquad\text{so}\qquad
i_{P}(C_{-},D)=2.
\]

\medskip

Finally, we can conclude that
\[
i_{P}(C,D)=i_{P}(C_{+},D)+i_{P}(C_{-},D)=2+2=4.
\]







\subsection*{Problem 7.12}

Let
\[
f(x,y)=(y-x^{2})(y+x^{2}), \qquad g(x,y)=y^{2}-x^{3}\in k[x,y],
\]
let $C:=V(f)$ and $D:=V(g)$ be plane curves in $\mathbb{A}^{2}$, and let $P=(0,0)$.
We need to compute $i_{P}(C,D)$.

\medskip

From Example 7.11. we have that
\[
i_{P}(C,D)=i_{P}(f,g)=i_{P}(f_{+},g)+i_{P}(f_{-},g)
= i_{P}(C_{+},D)+i_{P}(C_{-},D).
\]
Thus it suffices to compute $i_{P}(C_{+},D)$ and $i_{P}(C_{-},D)$.

\medskip

$P$ is nonsingular on each component. Since $P$ is nonsingular on $C_{+}$, Proposition 7.34 gives
\[
i_{P}(C_{+},D)=\operatorname{ord}^{C_{+}}_{P}(g).
\]

To compute this order, we first restrict $g$ to the curve $C_{+}$.
On $C_{+}$ we have $y=x^{2}$, so
\[
g|_{C_{+}} = g(x,x^{2}) = (x^{2})^{2}-x^{3}
= x^{4}-x^{3}
= x^{3}(x-1).
\]

Next, we choose a uniformizing parameter.
The tangent line to $C_{+}$ at $P$ is $y=0$ (this is computed in Example 7.6).
In particular, the line $x=0$ is \emph{not} tangent to $C_{+}$ at $P$.
Thus, by Theorem 7.23, $x$ is a uniformizing parameter for the DVR $\mathcal{O}_{C_{+},P}$, so
\[
\operatorname{ord}^{C_{+}}_{P}(x)=1.
\]

Finally, note that
\[
(x-1)\big|_{x=0}=-1\neq 0,
\]
so $x-1$ does not vanish at $P$ on $C_{+}$ and is therefore a unit in $\mathcal{O}_{C_{+},P}$ by Proposition 6.45.
Hence it contributes order $0$, and we obtain
\[
\operatorname{ord}^{C_{+}}_{P}\bigl(x^{3}(x-1)\bigr)
=\operatorname{ord}^{C_{+}}_{P}(x^{3})
=3\operatorname{ord}^{C_{+}}_{P}(x)=3.
\]
Therefore
\[
i_{P}(C_{+},D)=3.
\]

\medskip

Next, compute $i_{P}(C_{-},D)$ in the same way.
Again Proposition 7.34 gives
\[
i_{P}(C_{-},D)=\operatorname{ord}^{C_{-}}_{P}(g).
\]
On $C_{-}$ we have $y=-x^{2}$, so
\[
g|_{C_{-}} = g(x,-x^{2}) = (-x^{2})^{2}-x^{3}
= x^{4}-x^{3}
= x^{3}(x-1).
\]
As before, $P$ is a simple point of $C_{-}$ and its tangent line at $P$ is again $y=0$ (since the linear term is still $y$),
so $x=0$ is not tangent. Thus $x$ is a uniformizing parameter in $\mathcal{O}_{C_{-},P}$ by Theorem 7.23, and the factor
$x-1$ is a unit because it evaluates to $-1$ at $x=0$ (Proposition 6.45). Hence
\[
\operatorname{ord}^{C_{-}}_{P}(g)=3,
\qquad\text{so}\qquad
i_{P}(C_{-},D)=3.
\]

\medskip

Finally, we can conclude that
\[
i_{P}(C,D)=i_{P}(C_{+},D)+i_{P}(C_{-},D)=3+3=6.
\]


\end{document}

